\input header
\input ras
\footline={\hss\tenrm 9-\folio\hss}\pageno=1

\centerline{\bf GLOBAL FIXED POINT THEORY}\bigskip

\centerline{\bf \S9. The Holomorphic Fixed Point Formula}\medskip

First note the following.

{\QP{\bf 9.1. Lemma.} \it Every rational map \[f(z)\not\equiv z\]
of degree \[d\] has exactly \[d+1\]
fixed points, counted with multiplicity.\medskip}

Here, by definition, the {\bit multiplicity} of a fixed point
is equal to one whenever the multplier satisfies \[\lambda\ne 1\], and
is strictly greater than one otherwise. (Compare \S7.) As an example, in
the special case of a polynomial map of degree \[d\ge 2\],
there is exactly one fixed point at infinity, and hence
\[d\] finite fixed points counted with multiplicity.\smallskip

{\bf Proof.}
Conjugating \[f\] by a fractional linear automorphism if necessary, we
may assume that the point at infinity is not fixed by \[f\]. Hence we
can write \[f\] as a quotient \[f(z)=p(z)/q(z)\] where
the polynomial \[q(z)\] has degree equal to \[d\] and \[p(z)\] has
degree at most \[d\]. Now the equation
\[f(z)=z\] is equivalent to the polynomial
equation \[p(z)=zq(z)\] of degree \[d+1\], and the assertion follows.\QED
\smallskip

Both Fatou and Julia made use of a ``well known'' relation between the
multipliers at the fixed
points of a rational map. Consider first an isolated
fixed point \[f(z_0)=z_0\] of a holomorphic map in one complex variable.
If \[z_0\ne\infty\], we define the {\bit holomorphic index} of \[f\] at
\[z_0\] to be the residue
$$	\iota(f\,,\,z_0)\;=\;{1\over 2\pi i}\oint {dz\over z-f(z)}	$$
where we integrate in a small loop in the positive direction around \[z_0\].
As usual, if \[z_0\in\hat\C\] happens to be the point at infinity, then
we must first introduce the local uniformizing parameter \[\zeta=\phi(z)=1/z\].
In this case, we define the residue of \[f\] at \[\infty\] to be equal to
the residue of \[\phi\circ f\circ \phi^{-1}\] at the origin.

{\QP{\bf 9.2. Theorem.} {\it For any rational map \[\;f:\hat\C\to\hat\C\;\]
with \[f(z)\] not identically equal to \[z\], we have the relation
$$	\sum_{f(z)=z}\,\iota(f,z)\;=\; 1\,,	$$
to be summed over all fixed points.
In the case of a simple fixed point, where the multiplier
\[\lambda=f'(z_0)\] satisfies \[\lambda\ne 1\], the index is given by
$$	\iota(f,z_0)\;=\;{1\over 1-\lambda}\;.	$$
In any case, the index at \[z_0\] is a local analytic invariant. That
is, if \[\,g=\phi\circ f\circ\phi^{-1}\]
where \[\phi\] is a local holomorphic change of coordinate, then
\[\iota(f,z_0)=\iota(g,\phi(z_0))\].}\medskip}

\noindent
As an application, if \[f\] has one fixed point with multiplier very close
to 1, and hence with \[|\iota|\] large, then it must have at least one
other fixed point with \[|\iota|\] large and hence with \[\lambda\]
close to (or equal to) 1.\smallskip

{\bf Proof of 9.2.} Conjugating \[f\] by a linear fractional automorphism
if necessary, we may assume that the point at infinity is not fixed by \[f\].
Then \[f(z)\] remains bounded as \[|z|\to\infty\], and there is
a Laurent series expansion of the form
$$	\big( z-f(z)\big)^{-1}\;=\; z^{-1}+c_2z^{-2}+c_3z^{-3}+\cdots	$$
for large \[|z|\]. It follows that
the integral of \[\;{1\over 2\pi i}\cdot {dz\over z-f(z)}\;\] in a large loop
around the origin is equal to \[+1\]. Evidently this integral is equal to
the sum of the residues \[\iota(f,z_j)\] at the fixed points of \[f\];
hence the summation formula. The
computation of \[\iota(f,z_j)\] at a fixed point with multiplier \[\lambda_j
\ne 1\] is similar, using the Laurent series expansion of \[\big(z-f(z)
\big)^{-1}\] in a neighborhood of \[z_j\].

The proof that \[f\] is a local analytic
invariant, even at a multiple fixed point where \[\lambda=1\],
can be sketched as follows. Choose a 1-parameter family of perturbations
\[f_t\] of the given map \[f_0\]
so that \[f_t\] has distinct fixed points, all with multiplier
different from 1, for all small \[t\ne 0\]. For example we can set \[f_t(z) 
=f(z)+t\]. Thus a fixed point \[z_0\]
of \[f=f_0\] will split up into a cluster of nearby simple fixed points for
\[t\ne 0\].  Since the integral in a fixed loop around \[z_0\] varies
continuously with \[t\], and since this integral is a sum of residues which
are evidently local analytic invariants when \[t\ne 0\], it follows that
\[\iota(f_0\,,\,z_0)\] is also a local analytic invariant.\QED\smallskip

For generalizations of this formula, the reader is referred to
Atiyah and Bott.\medskip

{\bf Examples.} A rational map \[f(z)=c\] of degree zero has just one fixed
point, with multiplier zero and hence with index \[\iota(f,c)=1\]. A
rational map of degree one
usually has two distinct fixed points, and the relation
$$	{1\over 1-\lambda_1}+{1\over 1-\lambda_2}\;=\;1	$$
simplifies to \[\lambda_1\lambda_2=1\]. Any polynomial map \[p(z)\]
of degree two or more has a ``superattractive'' fixed point at infinity, with
multiplier zero and hence with index \[\iota(p,\infty)=1\]. Thus
the sum of the indices
of the {\it finite} fixed points of a polynomial map of degree \[\ge 2\]
is always zero.
For a polynomial of degree exactly two, the relation
$$      {1\over 1-\lambda_1}+{1\over 1-\lambda_2}\;=\;0 $$
for the finite fixed points simplifies to \[\;{1\over 2}
(\lambda_1+\lambda_2)=1\,\].\medskip

{\QP{\bf 9.3. Lemma.} {\it A fixed point with multiplier \[\lambda\ne 1\]
is attracting if and only if its index \[\iota\]
has real part \[{\cal R}(\iota)>{1\over 2}\].}\bigskip}

Geometrically, this is proved by noting that inversion carries the disk
\[1+D\] having the origin as boundary point to the appropriate half-plane.
Computationally, it can be proved by noting that \[{1\over 2}<{\cal R}
({1\over 1-\lambda})\] if and only if
$$	1\;<\;{1\over 1-\lambda}+{1\over 1-\bar\lambda}\;.	$$
Clearing denominators, this reduces easily to the required inequality
\[\lambda\bar\lambda<1\].\QED

One important consequence is the following.\medskip

{\QP{\bf 9.4. Corollary.} {\it Every rational map of degree two or more
must have either a repelling fixed point, or a fixed point
with \[\lambda= 1\], or both.}\bigskip}

{\bf Proof.} If the \[d+1\] fixed points were distinct and were all
attracting or neutral, then each index would have real part \[{\cal R}(\iota)
\ge{1\over 2}\], hence the sum would have real part greater than or equal to
\[{d+1\over 2}>1\]; but this would contradict the Fixed Point Formula.\QED

Since repelling points and parabolic points both belong to the Julia set,
this yields another proof of the following. (Compare \S4.4, as well as
\S4.3 and \S7.5.)

{\QP{\bf 9.5. Corollary.} \it The Julia set, for any rational map of
degree two or more, is always non-vacuous.\smallskip}
\lln\medskip

{\bf Problem 9-1.} If \[\;f(z)=z+\alpha z^2+\beta z^3+({\rm higher\;terms})
\,\],
with \[\alpha\ne 0\], show that the holomorphic index is given by
\[\iota(f,0)=\beta/\alpha^2\]. As an example,
consider the one-parameter family of cubic maps
$$	f_\alpha(z)\;=\;z+\alpha z^2+z^3	$$
with a double fixed point at the origin. Show that the remaining finite
fixed point of \[f_\alpha\]
is attracting if and only if \[\alpha^2\] lies within a unit disk centered
at \[-1\], or if and only if \[\alpha\] lies within a figure 8 shaped region
bounded by a lemniscate. Show that \[f_\alpha\] can be perturbed so that the
double fixed point at the origin splits up into two fixed points which
are {\it both} attractive
if and only if \[\alpha^2\] lies inside the disk of radius \[1/2\] centered
at \[1/2\], or if and only if \[\alpha\] lies within a region bounded
by a lemniscate shaped like the symbol \[\infty\].\smallskip

{\bf Problem 9-2.} Any fixed point \[z_0\] for \[f\] is evidently also a
fixed point for \[f^{\circ n}\]. If \[z_0\] is attracting [or repelling],
show that \[\iota(f^{\circ n},z_0)\] tends to the limit 1 [or 0] as
\[n\to\infty\]. If \[\;f(z)=z+\alpha z^k+({\rm higher\;terms})\,\], show that
\[\iota(f^{\circ n},0)\] tends to the limit \[k/2\].

\vfil\eject % \bigskip\bigskip
\footline={\hss\tenrm 10-\folio\hss}\pageno=1

\centerline{\bf \S10. Most Periodic Orbits Repel.}\bigskip

This section will prove the following theorem of Fatou. By a {\bit cycle}
we will mean simply a periodic orbit of \[f\]. Recall that a cycle is
called {\bit attracting}, {\bit neutral}, or {\bit repelling} according as its
multiplier \[\lambda\] satisfies \[|\lambda|<1\,,\;|\lambda|=1\], or
\[|\lambda|>1\].

{\QP{\bf 10.1. Theorem.} {\it Let \[f:\C\to\C\] be a rational map
of degree two or more. Then \[f\] has at most a finite number of
cycles which are attracting or neutral.}\medskip}

We will see in \S11 that there always exist infinitely many repelling
cycles. Shishikura has given
the sharp upper bound of \[2d-2\] for the number of attracting or neutral
cycles, using methods of quasi-conformal surgery. However the classical
proof, which is given here, shows only that this number is less than
or equal to \[6d-6\].\smallskip

First consider the case of an attracting cycle \[\,\hat z_0\mapsto \hat z_1\mapsto\cdots
\mapsto \hat z_m=\hat z_0\,\].
Each \[\hat z_j\] is an attracting fixed point for the \[m$-fold composition
\[f^{\circ m}\]. If \[\Omega_j\] is the immediate basin for \[\hat z_j\] under
\[f^{\circ m}\], then the union \[\Omega_0\cup\cdots\cup\Omega_{m-1}\] will be
called the {\bit immediate basin}$\,$ for our \[m$-cycle. It can be described
as the union of those components of the Fatou set which intersect the given
attracting cycle. We continue to assume that \[f\] has degree two or more.

{\QP{\bf 10.2. Lemma.} \it The immediate basin \[\Omega_0\cup\cdots\cup
\Omega_{m-1}\] for any attracting cycle
of \[f\] must contain at least one critical point of \[f\].\medskip}

{\bf Proof.} In the special case of an attracting fixed point, this was proved
in \S6.6. Applying this result to the \[m$-fold iterate of \[f\], we see
that the immediate basin \[\Omega_0\] for the fixed point \[\hat z_0\]
of \[f^{\circ m}\] must contain at least one critical point \[z_0\] of
\[f^{\circ m}\]. Let
\[\,z_0\mapsto z_1\mapsto\cdots\,\] be the orbit of this critical point
under \[f\]. By the chain rule,
at least one the the \[m\] points \[z_0\,,\,z_1\,,\,
\ldots\,,\,z_{m-1}\] must be a critical point for \[f\]. Since \[z_j
\in\Omega_j\], this proves the Lemma.\QED

Next consider a parabolic cycle \[\,\hat z_0\mapsto \hat z_1\mapsto\cdots
\mapsto \hat z_m=\hat z_0\,\], with multiplier
\[\lambda\] equal to a \[q$-th root of unity. Then the \[qm$-fold
iterate \[f^{\circ qm}\]
maps each \[\hat z_j\] to itself with multiplier \[\lambda^q\] equal to \[1\].
Recall that the number \[n\ge 1\] of attracting petals around \[\hat z_j\]
must be some multiple of \[q\]. (Compare 7.2 and 7.6.)\smallskip

{\bf Definition.} The union of
those components of the Fatou set which contain one of these \[n\] attracting
petals around one of the \[m\] points of the cycle will be called
the {\bit immediate basin} of this parabolic
cycle. Thus the number \[nm\] of connected components of
this immediate basin is some multiple of \[qm\].

{\QP{\bf 10.3. Lemma.} \it If \[f\] has degree \[d\ge 2\], then
the immediate basin for any parabolic cycle must also
contain at least one critical point.\medskip}

{\bf Proof.} In the case of a fixed point with multiplier \[\lambda=1\],
this was proved in \S7.10. The general case follows easily, just as in the
argument above.\QED

In fact, it is not difficult to check that there must be at least
\[n/q\] distinct critical points in such an immediate basin.
Combining 10.2 and 10.3 we obtain the following.

{\QP{\bf 10.4. Lemma.} {\it A rational map of degree \[d\ge 2\] can
have at most \[2d-2\] cycles which are attracting or parabolic.
Similarly, a polynomial map of degree \[d\ge 2\] can
have at most \[d-1\] cycles in the finite plane
which are attracting or parabolic.}\medskip}

{\bf Proof.} Since the immediate basins for distinct cycles are distinct
by definition, it follows from 10.2 and 10.3 that
the number of attracting or parabolic cycles is less
than or equal to the number of critical points. In the case of a polynomial
of degree \[d\], the number of finite critical points is clearly at
most \[d-1\]. In the case of a rational function of degree \[d\], by
the Riemann-Hurwitz formula \S5.1 the number of critical
ponts, counted with multiplicity, is equal to \[2d-2\]. Hence
the number of distinct critical points is at most \[2d-2\]. In either
case, the conclusion follows.\QED

{\bf 10.5. Remark.} This Lemma gives a practical computational procedure
for finding all attracting or parabolic cycles. We must simply follow the
orbits of all critical points, and test for convergence.

{\QP{\bf 10.6. Lemma.} {\it For a rational map of degree \[d\ge 2\],
the number of neutral cycles which have multiplier \[\lambda\ne 1\]
is at most \[4d-4\].}\medskip}

Evidently 10.4 and 10.6 together imply Theorem 10.1.
The proof of 10.6 begins as follows. Following Fatou, we perturb the given
map \[f\] and then applying 10.4. It will be convenient
to compare \[f\] with the map \[z\mapsto z^d\] of the same degree.
Note that this model map has no neutral cycles. It has
superattractive fixed points at zero and infinity,
but otherwise all of its periodic points are strictly repelling, with
\[|\lambda|=d^m>1\] where \[m\] is the period. Let
\[f(z)=p(z)/q(z)\] where \[p(z)\] and \[q(z)\] are polynomials, at least
one of which has degree \[d\]. Consider the one-parameter family of
rational maps
$$	f_t(z)\;= \;{(1-t)p(z)\,+\,tz^d\over (1-t)q(z)\,+\,t\;\;\;}	$$
with \[f_0(z)=f(z)\] and \[f_1(z)=z^d\]. There will be some finite set
of exceptional values of \[t\] for which the numerator and denominator
of this fraction have a common divisor. Algebraically this condition
is expressed by setting the ``resultant'' of the numerator and denominator
equal to zero.  Geometrically, it means that
a zero and pole of \[f_t\] crash together. If we exclude this finite set of
bad values of \[t\], then clearly \[f_t(z)\] is holomorphic as a function
of two variables, and each \[f_t\] has degree \[d\].

Suppose that \[f=f_0\] had \[4d-3\] distinct neutral cycles
with multipliers \[\lambda_j\ne 1\].
By the Implicit Function Theorem, we could follow each
of these cycles under a small deformation of \[f_0\]. Thus, for small
values of \[|t|\], the map \[f_t\] would have corresponding cycles
with multipliers \[\lambda_j(t)\] which depend holomorphically on \[t\],
with \[|\lambda_j(0)|=1\].

{\QP{\bf 10.7. Sub-Lemma.} {\it None of these functions \[t\mapsto
\lambda_j(t)\] can be constant throughout a neighborhood of \[t=0\].}\medskip}

{\bf Proof.} Suppose that for some \[j\] the function \[t\mapsto
\lambda_j(t)\] were constant throughout a neighborhood of \[t=0\]. Then
we will show that it is possible to continue analytically along any path
from 0 to 1 in the \[t$-plane which misses the finitely many
exceptional values of \[t\].
To prove this, we must check that the set of \[t\] for which we can continue
is both open and closed. But it is closed since any limit point of periodic
points with fixed multiplier \[\lambda_j\ne 1\] is itself a periodic point
with this same multiplier; and it is open since any such cycle varies
smoothly with \[t\], throughout some open neighborhood
in the \[t$-plane, by the Implicit Function Theorem.
Now continuing analytically to \[t=1\], we see that the map \[z\mapsto z^d\]
must also have  a cycle with mutiplier equal to \[\lambda_j\], with \[|
\lambda_j|=1\]. But this is known to be false, which proves 10.7.\QED

The proof of 10.6 continues as follows. We can express each of our \[4d-3\]
multipliers as a locally convergent power series
$$	\lambda_j(t)/\lambda_j(0)\;=\;1+a_j t^{n(j)}+({\rm higher\;terms})\,,
	 $$
where \[a_j\ne 0\] and \[n(j)\ge 1\].
Hence \[\log|\lambda_j(t)|\] is equal to the real part of\break
\[\;a_j t^{n(j)}+
{\rm (higher\;terms)}\,\]. If we ignore the higher order terms, this means
that we can divide the \[t$-plane into \[n(j)\] sectors for which
\[|\lambda_j(t)|>1\] and \[n(j)\] congruent sectors for which
\[|\lambda_j(t)|<1\]. Taking account of the higher order terms, we have
the following statement. Note that \[{\rm sgn}(
\log|\lambda|)\] is equal to \[+1\] or \[-1\] according as \[|\lambda|>1\] or
\[|\lambda|<1\].

{\QP{\bf Assertion.} \it The step function
$$	\theta\;\mapsto\;\sigma_j(\theta)\;\;=\;\;\lim_{r\to 0}
	\;{\rm sgn}\,\log|\lambda_j(r e^{i\theta})|	$$
is well defined with value \[\pm 1\]
except at finitely many jump discontinuities; and has average equal to
zero.\smallskip}

\noindent Therefore the sum \[\;\sigma_1(\theta)+\cdots+\sigma_{4d-3}
(\theta)\;\] is also a well defined step function with average zero.
Since this sum takes odd
values almost everywhere, we can choose some \[\theta\] for which
\[\sigma_1(\theta)+\cdots+\sigma_{4d-3}(\theta)\le -1\].
If we choose \[r\] sufficiently small and set \[t=re^{i\theta}\],
this means that \[f_t\] has at least \[2d-1\] distinct cycles
with multiplier satisfying \[|\lambda_j|<1\].
But this is impossible by 10.4, which proves 10.6 and 10.1.\QED

\vfil\eject % \bigskip\bigskip
\footline={\hss\tenrm 11-\folio\hss}\pageno=1

\centerline{\bf \S11. Repelling Cycles are Dense in \[J\].}\bigskip

We saw in \S4.3 that every repelling cycle is contained in the
Julia set. The following much sharper statement was proved
by Fatou and by Julia. (Compare 11.9.)

{\QP{\bf 11.1. Theorem.} {\it The Julia set for any rational map
of degree \[\ge 2\] is equal
to the closure of its set of repelling periodic points.}\medskip}

\noindent Since the proofs by Julia and by Fatou are interesting and
different, we will give both.\medskip

{\bf Proof following Julia.} We will use
the Holomorphic Fixed Point Formula of \S9.
Recall from 9.4 that every rational map \[f\] of degree two or more has either
a repelling fixed point, or a fixed point with \[\lambda=1\]. In either
case, this fixed point belongs to the Julia set \[J(f)\].
(Compare \S4.3 and \S7.5.)\smallskip

Thus we can start with a fixed point \[z_0\] in the Julia set. Let \[U
\subset\hat\C\] be any open set, disjoint from \[z_0\],
which intersects \[J(f)\]. The next step
is to construct a special orbit \[\;\cdots\mapsto z_2\mapsto z_1
\mapsto z_0\;\] which passes through \[U\] and terminates at this fixed point
\[z_0\]. By definition, such an orbit is called {\bit homoclinic} if
the backwards limit \[\lim_{j\to\infty} \;z_j\] exists and is equal to
the terminal point \[z_0\]. To construct a homoclinic orbit, we will
appeal to Theorem 4.8 which says that there exists an integer \[r>0\] and
a point \[z_r\in J(f)\cap U\] so that the \[r$-th forward image \[f^{\circ
r}(z_r)\] is equal to \[z_0\]. Given any neighborhood \[N_0\] of \[z_0\], we
can repeat this argument and conclude that there exists an integer \[q>r\]
and a point \[z_q\in N_0\] so that \[f^{\circ(q-r)}(z_q)=z_r\]. (Figure 11.)
\vskip .5in
\midinsert
\insertRaster homoc.ps pixels 480 by 640 scaled 500
\vskip -.3in
\centerline{Figure 11. A homoclinic orbit.}
\vfil
\endinsert

To be more explicit, in the case where \[z_0\] is a repelling fixed point
we choose \[N_0\] to be a linearizing neighborhood, as in the
Koenigs Theorem 6.1. In the parabolic case, we choose \[N_0\]
to be a flower neighborhood, as in 7.2. In either case, we choose \[N_0\]
small enough to be disjoint from \[z_r\]. It then follows that we can
inductively choose preimages \[\cdots\mapsto
z_j\mapsto z_{j-1}\mapsto\cdots\mapsto
z_q\], all inside of the neighborhood \[N_0\]. {\it These preimages \[z_j\]
will automatically converge to \[z_0\] as \[j\to\infty\].} If \[z_0\] is
repelling, this is clear. In the parabolic case, \[z_q\]
cannot belong to an attracting petal; hence it must belong to a repelling
petal, and again this statement is clear.\smallskip

First suppose that none of the points \[\;\cdots\mapsto z_j\mapsto\cdots
\mapsto z_0\;\] in this homoclinic orbit are critical points of \[f\].
Then a sufficiently small disk neighborhood \[V_q\] of\break
\[z_q\in N_0\] will map
diffeomorphically under \[f^{\circ q}\] onto a neighborhood \[V_0\] of
\[z_0\]. Pulling this neighborhood \[V_q\] back under iterates of \[f^{-1}\],
we obtain neighborhoods \[z_j\in V_j\] for all \[j\], shrinking down towards
the limit point \[z_0\] as \[j\to\infty\]. In particular, if we choose
\[p\] sufficiently large, then \[\bar V_p\subset V_0\]. Now \[f^{-p}\] maps
the simply connected open set \[V_0\] holomorphically into a compact subset
of itself. Hence it contracts the Poincar\'e metric on \[V_0\] by a
factor \[c<1\], and therefore must have a attractive fixed point \[z'\]
within \[V_p\]. Evidently this point \[z'\in V_p\] is a repelling periodic
point of period \[p\] under the map \[f\]. Since the orbit of \[z'\]
under \[f\] intersects the
required open set \[U\], the conclusion follows.
\bigskip

If our homoclinic orbit contains critical points, then this argument must
be modified very slightly as follows. We can still choose simply
connected neighborhoods \[V_j\] of the \[z_j\] so that \[\bar V_p\subset
V_0\] for some large \[p\], and so that \[f\] maps each \[V_j\] onto
\[V_{j-1}\]. However, some finite number of these mappings will be branched.
Choose a slit \[S\] in \[V_0\] from the boundary to the midpoint \[z_0\]
so as to be disjoint from
\[\bar V_p\], and choose some sector in \[V_p\] which maps isomorphically
onto \[V_0-S\] under \[f^{\circ p}\]. The proof now proceeds just as
before.\QED\medskip
% \advancepageno

{\bf Proof of 11.1 following Fatou.} In this case, the main idea is an easy
application of Montel's Theorem (\S2.5). However, we must use Theorem 10.1 to
finish the argument.\smallskip

To begin the proof, recall from 4.9 that the Julia set \[J(f)\] has
no isolated points. Hence we can exclude finitely many points of \[J(f)\]
without affecting the argument. Let \[z_0\] be any point of \[J(f)\]
which is not a fixed point, and not a critical value. In other words, we
assume that there are \[d\] preimages \[z_1\,,\,\ldots\,,\,z_d\], which
are distinct from each other and from \[z_0\], where \[d\ge 2\] is the degree.
By the Inverse Function Theorem, we can find \[d\] holomorphic functions
\[z\mapsto \varphi_j(z)\] which are defined throughout some neighborhood \[N\] of
\[z_0\], and which satisfy \[f(\varphi_j(z))=z\], with \[\varphi_j(z_0)=z_j\].
We claim that for some \[n>0\] and for some \[z\in N\] the function
\[f^{\circ n}(z)\] must take one of the three values \[z\,,\,
\varphi_1(z)\] or \[\varphi_2(z)\]. For otherwise the family of holomorphic
functions
$$	g_n(z)\;=\;	{\bigl(f^{\circ n}(z)-\varphi_1(z)\bigr)\,\bigl(
	z-\varphi_2(z)\bigr)
	\over \bigl(f^{\circ n}(z) -\varphi_2(z)\bigr)\,\bigl
	(z-\varphi_1(z)\bigr)}	$$
on \[N\] would avoid the three values \[0\,,\,1\] and \[\infty\], and hence
be a normal family. It would then follow easily that \[\{f^{\circ n}|N\}\]
was also a normal family, contradicting the hypothesis that \[N\]
intersects the Julia set. Thus we can find \[z\in N\] so as to satisfy either
\[f^{\circ n}(z)=z\] or \[f^{\circ n}(z)=\varphi_j(z)\]. Clearly it
follows that \[z\] is a periodic point of period \[n\] or \[n+1\]
respectively.\smallskip

This shows that every point in \[J(f)\] can be approximated arbitrarily
closely by periodic points. Since all but finitely many of these periodic
points must repel, this completes the proof.\QED
\smallskip

There are a number of interesting corollaries.

{\QP{\bf 11.2. Corollary.} {\it If \[U\] is an open set which intersects the
Julia set \[J\] of \[f\], then for \[n\] sufficiently large
the image \[\;f^{\circ n}(U\cap J)\;\] is equal
to the entire Julia\break set \[J\].}\medskip}

{\bf Proof.} We know that \[U\] contains a repelling periodic point \[z_0\] of
period say \[p\]. Thus \[z_0\] is fixed by the iterate \[g=f^{\circ p}\].
Choose a
small neighborhood \[V\subset U\] of \[z_0\] with the property that
\[V\subset g(V)\]. Then clearly \[V\subset g(V)\subset g^{\circ 2}
(V)\subset\cdots\,\]. But it follows from 4.6 or 4.8 that the union of
the open sets \[g^{\circ n}(V)\] contains the entire Julia set \[J=
J(f)=J(g)\].
Since \[J\] is compact, this implies that \[J\subset g^{\circ n}(V)\subset
g^{\circ n}(U)\] for \[n\] sufficiently large, and the corresponding
statement for \[f\] follows.\QED


{\QP{\bf 11.3. Corollary.} {\it If a Julia set \[J\] is not connected,
then it has uncountably many distinct connected components.}\medskip}

{\bf Proof.} Suppose that \[J\] is the union \[J_0\cup J_1\] of two disjoint
non-vacuous compact subsets. After replacing \[f\] by some iterate
\[g=f^{\circ n}\], we may assume by 11.2 that\break
\[g(J_0)=J\] and \[g(J_1)=J\].
Now to each point \[z\in J\] we can assign an infinite sequence of symbols
$$	\epsilon_0(z)\,,\,\epsilon_1(z)\,,\,\epsilon_2(z)\,,\,\ldots\;\;\in\;
	\{0,1\}		$$
by setting
\[\epsilon_k(z)\] equal to zero or one according as \[g^{\circ k}(z)\]
belongs to \[J_0\] or \[J_1\]. It is not difficult to check that
points with different symbol sequences must belong to different connected
components of \[J\], and that all possible symbol sequences actually occur.
\QED

{\bf Remark.} Evidently only countably many of the connected components of
\[J\] can have positive Lebesgue measure. I don't know whether
all but countably many of the components of \[J\] must be single points.
\smallskip

Siegel disks and Cremer points do not seem directly related to critical
points, yet there is a connection as follows. By a {\bit critical orbit}
we will mean the forward orbit of some critical point.

{\QP\bf 11.4. Corollary. {\it Every boundary point of a maximal
Siegel disk \[U\]
belongs to the closure of some critical orbit.}\medskip}

{\bf Proof.} Otherwise, we could construct a small disk \[V\] around
the given point \[z_0\in\partial U\] so that the forward orbits of all
critical points avoid \[V\]. This would mean that
every branch of the \[n$-fold iterated inverse function \[f^{-n}\]
could be defined as a single valued holomorphic function \[f^{-n}:V\to
\hat\C\]. Let us choose that particular branch which carries the intersection
\[U\cap V\] into \[U\] by a
``rotation'' of \[U\]. Since the rotation number is irrational, we can choose
some subsequence of these inverse maps which converges to the identity
map on \[U\cap V\]. This is evidently a normal family, since it avoids
the central
part of \[U\]. Hence there is a sub-sub-sequence \[\{f^{-n(i)}\}\] which
converges on all of \[V\], necessarily to the identity map of \[V\].
It follows easily that the corresponding
sequence of forward iterates \[f^{\circ n(i)}\] also converges to
the identity on \[V\]. But this contradicts 11.2.\QED

{\QP\bf 11.5. Corollary. {\it Every Cremer point
is a non-isolated point in the closure of some critical orbit.}\medskip}

{\bf Proof.} (This proof applies also to parabolic cycles. Compare \S10.3.) 
Otherwise we could choose a small disk \[V\]
around the given point \[z_0\] so that no critical orbit intersects the
punctured disk \[V\ssm\{z_0\}\]. Replacing \[f\] by some iterate if necessary,
we may assume that \[z_0\] is fixed by \[f\].
Arguing as above, there exists a unique
holomorphic branch  \[f^{-n}:V\to\hat\C\] of the \[n$-fold iterated inverse
function which fixes the point \[z_0\]. These inverse maps form a normal
family since, for example, they avoid any periodic orbit which is disjoint
from \[V\]. Thus we can choose a subsequence \[f^{-n(i)}\]
converging locally uniformly to some holomorphic map \[h:V\to\hat\C\].
By the Inverse function Theorem, since \[|h'(z_0)|=1\], this \[h\] must map some
small neighborhood of \[z_0\] isomorphically onto a neighborhood \[V'\] of \[z_0\].
It follows that the corresponding forward maps \[f^{\circ n(i)}\] converge
on \[V'\] to the inverse map \[h^{-1}:V'\to V\]. Again,
this contradicts 11.2.\QED
\smallskip

By definition, the rational map \[f\] is called {\bit post critically finite~}
(or is called a {\bit Thurston map\/}) if
it has the property that every critical orbit is finite, or in other words
is either periodic or eventually periodic. According to Thurston, such a
map can be uniquely specified by a finite topological description.
(Compare Douady \& Hubbard [DH1].)

{\QP\bf 11.6. Corollary. {\it If \[f\] is post critically finite,
then every periodic
orbit of \[f\] is either repelling or superattracting.}\smallskip}

{\QP\bf 11.7. Corollary. {\it More generally, suppose that \[f\] has the
property that every critical orbit either is finite, or converges to an
attracting periodic orbit. Then every periodic orbit of \[f\] is either
attracting or repelling; there are no parabolic cycles, Cremer cycles, 
or Siegel cycles.}\medskip}

\noindent (Compare 13.5, 14.4)
The proofs are immediate, and will be left to the reader.\QED

By definition, \[z_0\] belongs to the Julia set if {\it some} sequence of
iterates of \[f\] has no subsequence which converges throughout a neighborhood
of \[z_0\]. However, a priori, there could be other sequences which do
converge.

{\QP{\bf 11.8. Corollary.} \it If \[z_0\] belongs to the Julia set, then
no sequence of iterates of \[f\] can converge uniformly
throughout a neighborhood of \[z_0\].\medskip}

For if \[\{f^{\circ n(i)}\}\] converges uniformly,
with \[n(i)\to\infty\], then according to Weierstrass (\S1.4) the
sequence of derivatives \[df^{\circ n(i)}(z)/dz\] must also converge. But
if \[z_1\] is a repelling periodic point, then evidently the sequence of
derivatives at \[z_1\] diverges to infinity.\QED
\medskip

{\bf 11.9. Concluding Remark.} The statement that the Julia set is equal to the
closure of the set of repelling periodic points is actually true for an
arbitrary holomorphic map of an arbitrary Riemann surface, providing that we
exclude just one exceptional case. For transcendental functions this
was proved by Baker [Ba1], and for maps of a torus or a Hyperbolic surface
it follows easily from \S4. The unique exceptional case occurs for degree one
maps of \[\hat\C\]
which have just one parabolic fixed point --- for example the map \[f(z)=z+1\]
with \[J(f)=\{\infty\}\].

On the other hand, for a holomorphic map of a complex 1-dimensional manifold
with two or more connected components (for example where both components
map into one), the Julia set can clearly be much bigger than the repelling
orbit closure.

\vfill
\bigskip

%\rightline{jwm, version of 9-3-91}

\end
